\section{Related work}
\paragraph{Social Intelligence in LLMs}

Large language models (LLMs) have led to new technologies that manage to handle common social use cases, including voice assistants, email autocomplete \citep{gmail}, AI-assisted counseling \citep{aicounseling}, etc.
However, human social interactions are more complicated and diverse than those restricted uses,
exposing model limitations in extended contexts. \citet{sap2023neural} study the limitations of social intelligence in current LLMs and conclude that current models struggle with Theory of Mind tasks such as SocialIQa~\citep{sap-etal-2019-social} and ToMi~\citep{le-etal-2019-revisiting}. In the Avalon game setting, 
\citet{light2023text} show that it is still challenging for LLM agents to successfully deceive, deduce, and negotiate with other players, particularly in a multi-agent environment. A potential method to improve the Theory of Mind in language agents is through meta learning the mental model of the interlocutor \citep{zhu2021few,zhu2022language,liu2022computational}. We leave explicity modeling Theory of Mind in language agents to improve the social intelligence as the future work. These studies show that the effective development of general social intelligence in model training has yet to be fully realized.


Studies have looked into the potential of behavior cloning from observational data for enhancing social intelligence via interaction 
\citep{wang2023aligning}. 
\model echos social science theories of inferential social learning~\citep{inferential-social-learning}, where models learn not only by imitating but also by making inferences about social contexts.


\paragraph{Reinforcement Learning for LLMs}
\label{sec:rl_llm}
Reinforcement learning from human feedback (RLHF; \citet{christiano2017deep}) improves the alignment of LLMs to human preferences \citep{ouyang2022training}. Direct Preference Optimization~\cite{dpo} and $\Psi$ Policy Optimization \cite{ipo} improve RLHF by optimizing the LLM policy without relying on the reward model. These online RL methods often require online data collection, which has a longer latency in multi-agent settings. 

Typical types of offline self-reinforcement include self-imitation learning (SIL; \citet{oh2018self}), reward ranked fine-tuning (RAFT; \citet{dong2023raft}), and reinforced self-training (ReST; \citet{ReST}). SIL sets a replay buffer and imitates state-action pairs when it is better than the current value estimation. RAFT generates multiple outputs and utilizes the reward model to filter out a subset. ReST is a more complicated version of RAFT with multiple improve steps. \model applies offline self-reinforcement to training LLMs on social tasks and utilizes the GPT-4 to provide rewards for multi-turn social interaction. We leave investigating the effects of different offline methods on training social intelligence to future work.

\paragraph{LLM Alignment and Evaluation}
Advances in fine-tuning methods like parameter-efficient fine-tuning~\citep{li2021prefixtuning, lester2021power,lora} have  
improved LLMs' capabilities to better understand the restriction and rules given by human, enhancing their capability for social learning and interaction. Other governance objectives align LLM behaviors via robustness, interpretability, controllability, and ethicality
\citep{ji2024ai}. We focus on evaluating our trained LLMs' alignment with human social norms via safety and toxicity. 

It has been pointed out that continual fine-tuning can lead to catastrophic forgetting of LLMs, in terms of domain knowledge, reasoning, and reading comprehension~\citep{catastrophic-forgetting}. To test the general question answering and reasoning capabilities of our trained LLMs, we measure their performance on the Massive Multitask Language Understanding (MMLU) benchmark~\cite{hendrycks2020measuring}, a holistic benchmark designed to test the knowledge of a model across 57 subjects. 
